\section{Societal and Environmental Impact (M7)}
\label{sec:impact}

\subsection{Positive Potential}

SAR drones reduce search time by covering terrain faster than ground teams and detecting casualties in areas too dangerous or remote for human access~\cite{sar_drones_review,goodrich2008sar}. What makes \emph{our} system's contribution specific is the cost barrier it removes: a \pounds60 Raspberry Pi running a 3.2\,MB TFLite model makes autonomous visual search viable for volunteer mountain rescue teams, community first responders, and organisations in countries where helicopter searches are unaffordable. The browser-based ground station requires no specialist software---any phone or tablet with a browser becomes a command interface---which matters when the operator may be a volunteer, not an engineer. I designed these choices deliberately, not just for technical elegance but because I believe the gap between ``technology exists'' and ``technology is accessible'' is where engineering effort is most needed.

\subsection{Environmental Impact}

I made deliberate choices to minimise the project's environmental footprint, and I can quantify several of them. Simulation-first development replaced approximately 50 physical test flights; assuming each flight consumes a 6S LiPo cycle (${\sim}$2.5\,kg CO\textsubscript{2} from battery manufacturing amortisation, motor wear, and travel to the test site), this avoided an estimated 125\,kg of CO\textsubscript{2} equivalent. The Pi~5 consumes ${\sim}$15\,W under inference load versus 300\,W+ for a GPU server running the same model via cloud API; over a 20-minute mission, that is 5\,Wh versus 100\,Wh. Model training used Google Colab's free tier (shared GPU, ${\sim}$0.5\,kWh for 150 epochs) rather than a dedicated server. YOLOv8n was chosen as the smallest viable architecture, and automatic SSSI geofencing protects the site's rare bird habitats by switching to manual mode near the no-fly zone rather than relying on pilot awareness. The drone platform is reusable for future cohorts, and the software is released under MIT licence, avoiding the need to rebuild from scratch.

\subsection{Dual-Use and Adverse Impact}

The same technology that searches for a missing person can surveil one. I am not able to discuss this abstractly: as a Ukrainian student, the military application of autonomous surveillance drones is not a hypothetical scenario for me. I have seen what these systems do when the ``target'' is not a casualty to rescue but a person to track. This personal experience made me deliberate about every design choice that touches autonomy.

The system I built detects and localises, but a human operator decides what to do---human-in-the-loop verification is not a feature I added for the rubric, but a constraint I insisted on because I understand what happens when it is removed. Images are processed locally and never transmitted to a cloud service; detection logs are stored on the Pi's SD card with no persistent surveillance capability. The confidence threshold (0.4) was set conservatively to reduce false positives, because in SAR, a false detection wastes search time that a real casualty cannot afford.

However, I must be honest about the limits of these safeguards. The human-in-the-loop is a software constraint, not a hardware one---removing the verification prompt is a single-line code change. The MIT licence means anyone can fork the repository and strip the geofence, the verification stage, and the logging. The edge AI architecture I chose specifically to avoid cloud dependency also means the system operates with no external oversight. These are not theoretical risks; they are design properties of the system I built. If I were to commercialise this work, I would explore hardware-enforced geofencing (e.g., a GPS-locked relay on the power bus) and a code-signing mechanism that prevents modified firmware from running---measures that move safety from software policy to physical constraint~\cite{leveson2011}.
